\chapter{2010年重庆卷（文）}

\section{单选题}

\begin{problem}
\((x + 1)^{4}\) 的展开式中 \(x^{2}\) 的系数为（\quad）

\choices
    {\(4\)}
    {\(6\)}
    {\(10\)}
    {\(20\)}
\end{problem}

\begin{answer}
B
\end{answer}

\begin{solution}
由二项式定理，
\[
(x+1)^4=\mathrm C_4^0x^4+\mathrm C_4^1x^3+\mathrm C_4^2x^2+\mathrm C_4^3x+\mathrm C_4^4.
\]
因此 \(x^2\) 的系数为 \(\mathrm C_4^2=6\)，故选 B.
\end{solution}

\begin{problem}
在等差数列 \(\{a_{n}\}\) 中，\(a_{1} + a_{9} = 10\)，则 \(a_{5}\) 的值为（\quad）

\choices
    {\(5\)}
    {\(6\)}
    {\(8\)}
    {\(10\)}
\end{problem}

\begin{answer}
5
\end{answer}

\begin{solution}
由等差数列的性质，\(a_{1}+a_{9}=2a_{5}\)，所以 \(2a_{5}=10\)，从而 \(a_{5}=5\).
\end{solution}

\begin{problem}
若向量 \(\bs{a} = (3, m)\)，\(\bs{b} = (2, -1)\)，\(\bs{a} \cdot \bs{b} = 0\)，则实数 \(m\) 的值为（\quad）

\choices
    {\(-\frac{3}{2}\)}
    {\( \frac{3}{2} \)}
    {\(2\)}
    {\(6\)}
\end{problem}

\begin{answer}
D
\end{answer}

\begin{solution}
由题意，\(\bs{a}\cdot\bs{b}=3\times2+m\times(-1)=6-m=0\)，所以 \(m=6\).
\end{solution}

\begin{problem}
函数 \( y = \sqrt{16 - 4^x} \) 的值域是（\quad）

\choices
    {\([0, +\infty)\)}
    {\([0,4]\)}
    {\([0,4)\)}
    {\((0,4)\)}
\end{problem}

\begin{answer}
C
\end{answer}

\begin{solution}
由根式有意义，得 \(16-4^{x}\geq0\)，即 \(x\leq2\). 当 \(x\leq2\) 时，\(0<4^{x}\leq16\)，因此 \(0\leq y<4\)，且 \(y=0\) 在 \(x=2\) 时取得，故值域为 \([0,4)\)，选 C.
\end{solution}

\begin{problem}
某单位有职工 750 人，其中青年职工 350 人，中年职工 250 人，老年职工 150 人. 为了了解该单位职工的健康情况，用分层抽样的方法从中抽取样本. 若样本中的青年职工为 \(7\) 人，则样本容量为（\quad）

\choices
    {\(7\)}
    {\(15\)}
    {\(25\)}
    {\(35\)}
\end{problem}

\begin{answer}
B
\end{answer}

\begin{solution}
分层抽样中各层抽取比例相同，因此抽取比例为 \(\frac{7}{350}=\frac{1}{50}\). 所以样本容量为 \(750\times\frac{1}{50}=15\)，选 B.
\end{solution}

\begin{problem}
下列函数中，周期为 \(\pi\)，且在 \(\left[\frac{\pi}{4}, \frac{\pi}{2}\right]\) 上为减函数的是（\quad）

\choices
    {\( y = \sin \left(2x + \frac{\pi}{2}\right) \)}
    {\( y = \cos \left(2x + \frac{\pi}{2}\right) \)}
    {\( y = \sin \left(x + \frac{\pi}{2}\right) \)}
    {\( y = \cos \left(x + \frac{\pi}{2}\right) \)}
\end{problem}

\begin{answer}
A
\end{answer}

\begin{solution}
函数 \(y=\sin\left(2x+\frac{\pi}{2}\right)=\cos2x\) 的周期为 \(\pi\). 当 \(x\in\left[\frac{\pi}{4},\frac{\pi}{2}\right]\) 时，\(2x\in\left[\frac{\pi}{2},\pi\right]\)，而 \(\cos2x\) 在该区间上为减函数，故 A 符合条件. 对 B，\(y=\cos\left(2x+\frac{\pi}{2}\right)=-\sin2x\)，由 \(y'=-2\cos2x\geqslant0\)可知其在该区间上为增函数；C，D 的最小正周期均为 \(2\pi\)，不满足周期条件，因此选 A.
\end{solution}

\begin{problem}
设变量 \(x\)，\(y\) 满足约束条件 \(\begin{cases}
x \geqslant 0,\\
x - y \geqslant 0,\\
2x - y - 2 \leqslant 0
\end{cases}\)，则 \(z = 3x - 2y\) 的最大值为（\quad）

\choices
    {\(0\)}
    {\(2\)}
    {\(4\)}
    {\(6\)}
\end{problem}

\begin{answer}
C
\end{answer}

\begin{solution}
约束条件等价于 \(x\geq0\)，\(y\leq x\)，\(y\geq2x-2\). 由 \(2x-2\leq x\) 得 \(0\leq x\leq2\). 对固定的 \(x\)，因 \(z=3x-2y\) 随 \(y\) 减小而增大，故取 \(y=2x-2\)，于是 \(z=3x-2(2x-2)=4-x\leq4\). 当 \((x,y)=(0,-2)\) 时等号成立，所以最大值为 \(4\)，选 C.
\end{solution}

\begin{problem}
若直线 \(y = x - b\) 与曲线 \(\begin{cases}
x = 2 + \cos \theta,\\
y = \sin \theta
\end{cases}\)（\(\theta \in [0, 2\pi)\)）有两个不同的公共点，则实数 \(b\) 的取值范围为（\quad）

\choices
    {\((2 - \sqrt{2}, 1)\)}
    {\([2 - \sqrt{2}, 2 + \sqrt{2}]\)}
    {\((-\infty, 2 - \sqrt{2}) \cup (2 + \sqrt{2}, +\infty)\)}
    {\((2 - \sqrt{2}, 2 + \sqrt{2})\)}
\end{problem}

\begin{answer}
D
\end{answer}

\begin{solution}
由参数方程，曲线是圆 \((x-2)^2+y^2=1\). 联立直线方程得 \(b=x-y=2+\cos\theta-\sin\theta=2+\sqrt{2}\cos\left(\theta+\frac{\pi}{4}\right)\). 因而 \(b\) 的取值范围为 \([2-\sqrt{2},2+\sqrt{2}]\). 当 \(b\) 取两个端点时直线与圆相切，只有一个公共点；要有两个不同的公共点，必须有 \(2-\sqrt{2}<b<2+\sqrt{2}\)，选 D.
\end{solution}

\begin{problem}
到两互相垂直的异面直线的距离相等的点（\quad）

\choices
    {只有 \(1\) 个}
    {恰有 \(3\) 个}
    {恰有 \(4\) 个}
    {有无穷多个}
\end{problem}

\begin{answer}
D
\end{answer}

\begin{solution}
设两直线的公垂线长为 \(h\)，其中 \(h>0\). 建立空间直角坐标系，使两直线分别为
\[
l_1:\ y=z=0,\qquad l_2:\ x=0,\ z=h.
\]
对任意点 \(P(x,y,z)\)，其到两直线的距离分别满足
\[
d(P,l_1)^2=y^2+z^2,\qquad d(P,l_2)^2=x^2+(z-h)^2.
\]
令 \(y=0\)，并令 \(z=\frac{x^2+h^2}{2h}\)，则
\[
d(P,l_1)^2=z^2=x^2+(z-h)^2=d(P,l_2)^2.
\]
随着 \(x\) 取不同实数，可得到无穷多个满足条件的点，因此选 D.
\end{solution}

\begin{problem}
某单位拟安排 6 位员工在今年 6 月 14 日至 16 日（端午节假期）值班，每天安排 2 人，每人值班 1 天. 若 6 位员工中的甲不值 14 日，乙不值 16 日，则不同的安排方法共有（\quad）

\choices
    {\(30\text{ 种}\)}
    {\(36\text{ 种}\)}
    {\(42\text{ 种}\)}
    {\(48\text{ 种}\)}
\end{problem}

\begin{answer}
C
\end{answer}

\begin{solution}
每天的两人不计先后，因此不受限制的安排总数为 \(\frac{6!}{(2!)^3}=90\). 设事件 \(A\) 为甲值 14 日，事件 \(B\) 为乙值 16 日. 若甲值 14 日，先从其余 5 人中选 1 人与甲同日，再从剩余 4 人中选 2 人值 15 日，故 \(|A|=\mathrm C_5^1\mathrm C_4^2=30\). 同理，\(|B|=30\).

若甲值 14 日且乙值 16 日，甲的同日同事有 \(4\) 种选法，乙的同日同事有 \(3\) 种选法，故 \(|A\cap B|=4\times3=12\). 由容斥原理，满足条件的安排数为
\[
90-30-30+12=42.
\]
因此选 C.
\end{solution}

\section{填空题}

\begin{problem}
设 \(A = \{x \mid x + 1 > 0\}\)，\(B = \{x \mid x < 0\}\)，则 \(A \cap B =\fillinblank{}\).
\end{problem}

\begin{answer}
\(\{x \mid -1 < x < 0\}\)
\end{answer}

\begin{solution}
由 \(x+1>0\) 得 \(x>-1\)，又由 \(x<0\)，所以 \(A\cap B=\{x\mid -1<x<0\}\).
\end{solution}

\begin{problem}
已知 \(t > 0\)，则函数 \(y = \frac{t^2 - 4t + 1}{t}\) 的最小值为\fillinblank{}.
\end{problem}

\begin{answer}
\(-2\)
\end{answer}

\begin{solution}
\(y=t+\frac{1}{t}-4\). 由于 \(t>0\)，由基本不等式得 \(t+\frac{1}{t}\geq2\sqrt{t\cdot\frac{1}{t}}=2\)，等号当且仅当 \(t=1\) 时成立，所以 \(y\geq-2\). 因此函数的最小值为 \(-2\).
\end{solution}

\begin{problem}
已知过抛物线 \( y^{2} = 4x \) 的焦点 \(F\) 的直线交该抛物线于 \(A\)，\(B\) 两点，\( |AF| = 2 \)，则 \( |BF| = \)\fillinblank{}.
\end{problem}

\begin{answer}
\(2\)
\end{answer}

\begin{solution}
抛物线 \(y^2=4x\) 的参数方程可写为 \((x,y)=(u^2,2u)\)，其焦点为 \(F(1,0)\). 设 \(A\)、\(B\) 对应的参数分别为 \(u_1\)、\(u_2\). 两点弦的方程为
\[
x-\frac{u_1+u_2}{2}y+u_1u_2=0.
\]
因为该弦经过焦点 \(F(1,0)\)，所以 \(1+u_1u_2=0\)，即 \(u_1u_2=-1\).

对参数为 \(u\) 的点，
\[
|PF|=\sqrt{(u^2-1)^2+(2u)^2}=\sqrt{(u^2+1)^2}=u^2+1.
\]
由 \(|AF|=u_1^2+1=2\) 得 \(u_1^2=1\)，从而由 \(u_1u_2=-1\) 得 \(u_2^2=1\). 因此 \(|BF|=u_2^2+1=2\).
\end{solution}

\begin{problem}
加工某一零件需经过三道工序，设第一，二，三道工序的次品率分别为 \(\frac{1}{70}\)，\(\frac{1}{69}\)，\(\frac{1}{68}\)，且各道工序互不影响，则加工出来的零件的次品率为\fillinblank{}.
\end{problem}

\begin{answer}
\(\frac{3}{70}\)
\end{answer}

\begin{solution}
三道工序都合格的概率为 \(\left(1-\frac{1}{70}\right)\left(1-\frac{1}{69}\right)\left(1-\frac{1}{68}\right)=\frac{69}{70}\cdot\frac{68}{69}\cdot\frac{67}{68}=\frac{67}{70}\). 因此加工出来的零件为次品的概率为 \(1-\frac{67}{70}=\frac{3}{70}\).
\end{solution}

\begin{problem}
如图，图中的实线是由三段圆弧连接而成的一条封闭曲线 \(C\)，各段弧所在的圆经过同一点 \(P\)（点 \(P\) 不在 \(C\) 上）且半径相等. 设第 \(i\) 段弧所对的圆心角 \(\alpha_i\)（\(i = 1,2,3\)），则 \(\cos \frac{\alpha_1}{3} \cos \frac{\alpha_2 + \alpha_3}{3} - \sin \frac{\alpha_1}{3} \sin \frac{\alpha_2 + \alpha_3}{3} =\)\fillinblank{}.

\FigureLayoutDeclare{img/2010/chongqing_liberal/q15_fig1.png}{8.0cm}{16bp 10bp 17bp 6bp}
\bitmapfigure[max width=0.38\linewidth]{img/2010/chongqing_liberal/q15_fig1.png}
\end{problem}

\begin{answer}
\(-\frac{1}{2}\)
\end{answer}

\begin{solution}
设三圆的圆心分别为 \(O_1\)，\(O_2\)，\(O_3\)，并记 \(\beta_1=\angle O_2PO_3\)，\(\beta_2=\angle O_3PO_1\)，\(\beta_3=\angle O_1PO_2\)，则
\[
\beta_1+\beta_2+\beta_3=2\pi.
\]
设第 \(i\) 个圆与另外两个圆的另一个交点分别为 \(Q_{ij}\) 和 \(Q_{ik}\)，其中 \(\{i,j,k\}=\{1,2,3\}\). 由于两圆半径相等且都经过 \(P\)，四边形 \(P O_i Q_{ij} O_j\) 是菱形. 因此，在以 \(O_i\) 为圆心的圆上，从 \(Q_{ij}\) 到 \(P\) 的圆心角为 \(\pi-\beta_k\)，从 \(P\) 到 \(Q_{ik}\) 的圆心角为 \(\pi-\beta_j\). 经过 \(P\) 的那段弧所对的圆心角为
\[
(\pi-\beta_k)+(\pi-\beta_j)=2\pi-(\beta_j+\beta_k).
\]
\(\alpha_i\) 是其余实线圆弧所对的圆心角，所以 \(\alpha_i=\beta_j+\beta_k\). 从而
\[
\alpha_1+\alpha_2+\alpha_3=2(\beta_1+\beta_2+\beta_3)=4\pi.
\]
利用余弦的和角公式，原式为
\[
\cos\left(\frac{\alpha_1+\alpha_2+\alpha_3}{3}\right)=\cos\frac{4\pi}{3}=-\frac{1}{2}.
\]
\end{solution}

\section{解答题}

\begin{problem}
已知 \(\{a_n\}\) 是首项为 19，公差为 \(-2\) 的等差数列，\(S_n\) 为 \(\{a_n\}\) 的前 \(n\) 项和.
\begin{enumerate}
\item 求通项 \( a_{n} \) 及 \( S_{n} \)；
\item 设 \( \{b_{n}-a_{n}\} \) 是首项为 1，公比为 3 的等比数列，求数列 \( \{b_{n}\} \) 的通项公式及前 \(n\) 项和 \( T_{n} \).
\end{enumerate}
\end{problem}

\begin{solution}
\begin{enumerate}
\item 等差数列的通项公式为 \(a_n=a_1+(n-1)d=19-2(n-1)=21-2n\). 因而 \(S_n=\frac{n(a_1+a_n)}{2}=\frac{n(19+21-2n)}{2}=n(20-n)=20n-n^2\).
\item 由题意，\(b_n-a_n=3^{n-1}\)，所以 \(b_n=21-2n+3^{n-1}\). 因而 \(T_n=\sum_{k=1}^{n}b_k=\sum_{k=1}^{n}(21-2k)+\sum_{k=1}^{n}3^{k-1}=S_n+\frac{3^n-1}{2}=20n-n^2+\frac{3^n-1}{2}\).
\end{enumerate}
\end{solution}

\begin{problem}
在甲、乙等 6 个单位参加的一次“唱读讲传”演出活动中，每个单位的节目集中安排在一起. 若采用抽签的方式随机确定各单位的演出顺序（序号为 \(1,2,\ldots,6\)），求：
\begin{enumerate}
\item 甲、乙两单位的演出序号均为偶数的概率；
\item 甲、乙两单位的演出序号不相邻的概率.
\end{enumerate}
\end{problem}

\begin{solution}
把 6 个单位的演出顺序看作 6 个不同元素的全排列，样本空间的基本事件总数为 \(6!\).
\begin{enumerate}
\item 偶数序号有 \(2\)，\(4\)，\(6\) 共 \(3\) 个. 先给甲、乙安排两个不同的偶数序号，有 \(3\times2\) 种方法；其余 4 个单位任意排列，有 \(4!\) 种方法. 因此所求概率为 \(\frac{3\times2\times4!}{6!}=\frac{1}{5}\).
\item 甲、乙相邻时，把甲、乙看作一个整体，该整体内部有甲乙、乙甲两种顺序，故相邻的排列数为 \(2\times5!\). 因此不相邻的概率为 \(1-\frac{2\times5!}{6!}=\frac{2}{3}\).
\end{enumerate}
\end{solution}

\begin{problem}
设 \(\triangle ABC\) 的内角 \(A\)，\(B\)，\(C\) 的对边长分别为 \(a\)，\(b\)，\(c\)，且 \(3b^{2} + 3c^{2} - 3a^{2} = 4\sqrt{2}bc\).
\begin{enumerate}
\item 求 \(\sin A\) 的值；
\item 求 \(\frac{2\sin\left(A + \frac{\pi}{4}\right)\sin\left(B + C + \frac{\pi}{4}\right)}{1 - \cos 2A}\) 的值.
\end{enumerate}
\end{problem}

\begin{solution}
\begin{enumerate}
\item 由余弦定理，\(a^2=b^2+c^2-2bc\cos A\). 代入题设得 \(3(b^2+c^2-a^2)=4\sqrt{2}bc\)，即 \(6bc\cos A=4\sqrt{2}bc\). 由于 \(b\) 与 \(c\) 均为正数，所以 \(\cos A=\frac{2\sqrt{2}}{3}\). 又 \(0<A<\pi\)，故 \(\sin A=\sqrt{1-\cos^2A}=\sqrt{1-\frac{8}{9}}=\frac{1}{3}\).
\item 因为 \(B+C=\pi-A\)，所以 \(\sin\left(B+C+\frac{\pi}{4}\right)=\sin\left(\frac{5\pi}{4}-A\right)=\frac{\sqrt{2}}{2}(\sin A-\cos A)\). 于是原式为
\[
\frac{2\cdot\frac{\sin A+\cos A}{\sqrt{2}}\cdot\frac{\sqrt{2}}{2}(\sin A-\cos A)}{1-\cos2A}
=\frac{\sin^2A-\cos^2A}{2\sin^2A}
=\frac{\frac{1}{9}-\frac{8}{9}}{\frac{2}{9}}
=-\frac{7}{2}.
\]
\end{enumerate}
\end{solution}

\begin{problem}
已知函数 \( f(x) = ax^3 + x^2 + bx \)（其中常数 \(a\)，\(b \in \R\)），\( g(x) = f(x) + f'(x) \) 是奇函数.
\begin{enumerate}
\item 求 \( f(x) \) 的表达式；
\item 讨论 \( g(x) \) 的单调性，并求 \( g(x) \) 在区间 \( [1,2] \) 上的最大值和最小值.
\end{enumerate}
\end{problem}

\begin{solution}
\begin{enumerate}
\item \(f'(x)=3ax^2+2x+b\)，所以 \(g(x)=ax^3+(1+3a)x^2+(b+2)x+b\). 因为 \(g(x)\) 是奇函数，偶次项系数必须为零，故 \(1+3a=0\)，\(b=0\). 因而 \(a=-\frac{1}{3}\)，\(f(x)=-\frac{1}{3}x^3+x^2\).
\item 此时 \(g(x)=-\frac{1}{3}x^3+2x\)，所以 \(g'(x)=2-x^2\). 当 \(x<-\sqrt{2}\) 或 \(x>\sqrt{2}\) 时，\(g'(x)<0\)；当 \(-\sqrt{2}<x<\sqrt{2}\) 时，\(g'(x)>0\). 因此 \(g(x)\) 在 \((-\infty,-\sqrt{2})\) 和 \((\sqrt{2},+\infty)\) 上单调递减，在 \((-\sqrt{2},\sqrt{2})\) 上单调递增.

在 \([1,2]\) 上，\(g(x)\) 先增后减，故最大值为 \(g(\sqrt{2})=2\sqrt{2}-\frac{(\sqrt{2})^3}{3}=\frac{4\sqrt{2}}{3}\). 又 \(g(1)=\frac{5}{3}\)，\(g(2)=\frac{4}{3}\)，所以最小值为 \(\frac{4}{3}\).
\end{enumerate}
\end{solution}

\begin{problem}
如图，四棱锥 \(P - ABCD\) 中，底面 \(ABCD\) 为矩形，\(PA \perp\) 底面 \(ABCD\)，\(PA = AB = \sqrt{2}\)，点 \(E\) 是棱 \(PB\) 的中点.

\begin{enumerate}
\item 证明：\(AE \perp\) 平面 \(PBC\)；
\item 若 \(AD = 1\)，求二面角 \(B - EC - D\) 的平面角的余弦值.
\end{enumerate}

\bitmapfigure[max width=0.42\linewidth]{img/2010/chongqing_liberal/q20_fig1.png}
\end{problem}

\begin{solution}
\begin{enumerate}
\item 建立空间直角坐标系，取 \(A\) 为原点，令 \(AB\)，\(AD\)，\(AP\) 分别为 \(x\)，\(y\)，\(z\) 轴正方向，则 \(A(0,0,0)\)，\(B(\sqrt{2},0,0)\)，\(D(0,1,0)\)，\(C(\sqrt{2},1,0)\)，\(P(0,0,\sqrt{2})\)，\(E\left(\frac{\sqrt{2}}{2},0,\frac{\sqrt{2}}{2}\right)\). 于是 \(\overrightarrow{AE}=\left(\frac{\sqrt{2}}{2},0,\frac{\sqrt{2}}{2}\right)\)，\(\overrightarrow{PB}=(\sqrt{2},0,-\sqrt{2})\)，\(\overrightarrow{BC}=(0,1,0)\). 因而 \(\overrightarrow{AE}\cdot\overrightarrow{PB}=0\)，\(\overrightarrow{AE}\cdot\overrightarrow{BC}=0\). 直线 \(PB\) 与 \(BC\) 相交且都在平面 \(PBC\) 内，所以 \(AE\perp\) 平面 \(PBC\).
\item 由 \(PA\perp\) 底面 \(ABCD\) 及矩形性质，\(BC\perp AB\) 且 \(BC\perp PA\)，从而 \(BC\perp\) 平面 \(PAB\)，所以 \(BC\perp BE\). 又 \(PB=\sqrt{PA^2+AB^2}=2\)，所以 \(BE=1\)，进而 \(CE=\sqrt{BE^2+BC^2}=\sqrt{2}\). 由 \(AD\perp\) 平面 \(PAB\) 及 \(AE=1\)，得 \(DE=\sqrt{AD^2+AE^2}=\sqrt{2}\)，而 \(CD=AB=\sqrt{2}\). 所以 \(\triangle CED\) 为等边三角形.

取 \(F\) 为 \(CE\) 的中点. 在等边三角形 \(CED\) 中，\(DF\perp CE\)，且 \(DF=\frac{\sqrt{6}}{2}\)；在等腰直角三角形 \(EBC\) 中，\(BF\perp CE\)，且 \(BF=\frac{\sqrt{2}}{2}\). 因此 \(\angle BFD\) 是二面角 \(B-EC-D\) 的平面角. 又 \(BD^2=AB^2+AD^2=3\)，由余弦定理得
\[
\cos\angle BFD=\frac{BF^2+DF^2-BD^2}{2BF\cdot DF}=\frac{\frac{1}{2}+\frac{3}{2}-3}{2\cdot\frac{\sqrt{2}}{2}\cdot\frac{\sqrt{6}}{2}}=-\frac{\sqrt{3}}{3}.
\]
\end{enumerate}
\end{solution}

\begin{problem}
已知以原点 \(O\) 为中心，\(F(\sqrt{5}, 0)\) 为右焦点的双曲线 \(C\) 的离心率 \(e = \frac{\sqrt{5}}{2}\).

\begin{enumerate}
\item 求双曲线 \(C\) 的标准方程及其渐近线方程；
\item 如图，已知过点 \(M(x_1, y_1)\) 的直线 \(l_1: x_1x + 4y_1y = 4\) 与过点 \(N(x_2, y_2)\) 的直线 \(l_2: x_2x + 4y_2y = 4\) 的交点 \(E\) 在双曲线 \(C\) 上，直线 \(MN\) 与两条渐近线分别交于 \(G\)，\(H\) 两点，求 \(\overrightarrow{OG} \cdot \overrightarrow{OH}\) 的值.
\end{enumerate}

\bitmapfigure[max width=0.42\linewidth]{img/2010/chongqing_liberal/q21_fig1.png}
\end{problem}

\begin{solution}
\begin{enumerate}
\item 设双曲线的实半轴长为 \(a\)，虚半轴长为 \(b\)，焦半距为 \(c\). 由右焦点坐标得 \(c=\sqrt{5}\)，由离心率得 \(e=\frac{c}{a}=\frac{\sqrt{5}}{2}\)，所以 \(a=2\). 又 \(c^2=a^2+b^2\)，故 \(b^2=5-4=1\). 因而双曲线方程为 \(\frac{x^2}{4}-y^2=1\)，渐近线方程为 \(y=\pm\frac{x}{2}\).
\item 由直线 \(l_1\) 经过 \(M(x_1,y_1)\)，直线 \(l_2\) 经过 \(N(x_2,y_2)\)，分别代入直线方程，得 \(x_1^2+4y_1^2=4\) 和 \(x_2^2+4y_2^2=4\)，所以 \(M\)，\(N\) 都在椭圆 \(x^2+4y^2=4\) 上. 设 \(E=(X,Y)\). 由于 \(E\) 在两条切线 \(l_1\)，\(l_2\) 上，有
\[
Xx_1+4Yy_1=4,\qquad Xx_2+4Yy_2=4.
\]
因此 \(M\)，\(N\) 都在直线 \(Xx+4Yy=4\) 上，故直线 \(MN\) 的方程为 \(Xx+4Yy=4\). 又点 \(E\) 在双曲线 \(C\) 上，所以 \(\frac{X^2}{4}-Y^2=1\)，即 \(X^2-4Y^2=4\).

直线 \(MN\) 与渐近线 \(y=\frac{x}{2}\) 的交点记为 \(G\)，与渐近线 \(y=-\frac{x}{2}\) 的交点记为 \(H\). 代入 \(Xx+4Yy=4\)，得 \(G=\left(\frac{4}{X+2Y},\frac{2}{X+2Y}\right)\)，\(H=\left(\frac{4}{X-2Y},-\frac{2}{X-2Y}\right)\). 因此
\[
\overrightarrow{OG}\cdot\overrightarrow{OH}
=\frac{16}{(X+2Y)(X-2Y)}-\frac{4}{(X+2Y)(X-2Y)}
=\frac{12}{X^2-4Y^2}=3.
\]
\end{enumerate}
\end{solution}
